\begin{studiumworked}{M28. Una maniobra cambia dirección sin cambiar rapidez · INTERMEDIATE}
\textbf{Enunciado original.} En un modelo impulsivo, dos velocidades tienen igual módulo 7500 m/s y forman un ángulo de 3 grados. Determine el módulo de su diferencia y explique por qué no es cero aunque no cambie la rapidez.

\textbf{Datos utilizados.} Rapidez común: $7500\,\mathrm{m/s}$; Ángulo: $0.05235988\,\mathrm{rad}$.

\textbf{Incógnitas.} Magnitudes solicitadas y cierre físico.

\textbf{Hipótesis.} Modelo académico con hipótesis del capítulo. Datos originales; no representan hardware ni ensayo.

\textbf{Modelo físico.} Conservación y relaciones constitutivas declaradas.
\[|\Delta\mathbf v|^2=v^2+v^2-2v^2\cos\Delta i\]
\[\Delta v=2v\sin(\Delta i/2)\]
\textbf{Despejes y sustituciones.}
\begin{enumerate}
\item Se construye la diferencia vectorial, no la diferencia de módulos. La ley del coseno da $2v^2(1-\cos\Delta i)$.
\item La identidad $1-\cos\Delta i=2\sin^2(\Delta i/2)$ permite extraer raíz positiva. Sustituir $3\pi/180$ rad da $392.6542\,\mathrm{m/s}$.
\item La aproximación de ángulo pequeño $v\Delta i$ da $392.6991\,\mathrm{m/s}$. La proximidad comprueba el orden de magnitud, pero la expresión trigonométrica es la del modelo impulsivo exacto.
\end{enumerate}
\textbf{Resultados.} $\Delta v=392.6542$ m/s.

\textbf{Interpretación.} El resultado corresponde a las hipótesis indicadas; su interpretación requiere conservar fronteras y unidades.

\textbf{Comprobación.} El resultado tiende a cero con ángulo cero y a 2v para velocidades opuestas.

\textbf{Límites.} No incluye duración, gravedad durante impulso ni control de actitud.
\end{studiumworked}
